\section{Problem Formulation}
\label{sec:problem-formulation}

\subsection{The Problem}

Crude oil is simultaneously a physical commodity, a financial asset, and a
geopolitical bargaining chip, and its price reflects the interaction of all
three. A single, static statistical model that assumes one fixed relationship
between oil prices and a handful of macroeconomic drivers struggles to cope
with this environment: the relevant set of drivers and the strength of their relationships both change over time as the market
moves between demand-driven, supply-driven, and purely speculative/financial
regimes \autocite{kilian2009oil,hamilton2009causes}. A risk manager who prices
energy exposure, hedges an oil-linked book, or sets capital against commodity
risk therefore faces two compounding problems: (i) identifying which set of financial, macroeconomic, and physical-market variables is actually
driving the price at a given point in time, out of a large and only loosely
related candidate set, and (ii) doing so in a way that is transparent enough
to audit.

This project follows the framing used \cite{alvi2018pgm}: treat crude oil price forecasting as a problem of structure discovery over a mixed panel of financial, macroeconomic, and geopolitical/physical-market series, rather than a problem of fitting one fixed regression equation. Given
$m$ candidate variables $X_1, \ldots, X_m$ (our data dictionary in Table~\ref{tab:data-dictionary} lists 24: 5 financial, 8 macroeconomic, and 11 geopolitical/physical-market series sourced from Yahoo Finance, FRED, and the EIA Short-Term Energy Outlook respectively), the task is to learn both which variables are conditionally dependent on which others, and the strength and direction of those dependencies, directly from the data.

\subsection{Why Probabilistic Graphical Models Fit}

A Bayesian (belief) network represents the joint distribution of
$X_1, \ldots, X_m$ as a directed acyclic graph (DAG) in which each node's distribution is conditioned only on its direct parents, factorising the joint density as
\[
  p(x_1, \ldots, x_m) = \prod_{v=1}^{m} p\bigl(x_v \mid \mathrm{pa}(x_v)\bigr).
\]
This representation is a natural fit for the oil-pricing problem for three
reasons.

\begin{enumerate}
  \item \textbf{It encodes conditional independence, not just correlation.}
    Many of our 24 candidate series are themselves highly correlated with
    each other (e.g.\ WTI and Brent futures, or the several EIA supply/demand
    series derived from the same underlying balance identity — see the
    correlation heatmap in Figure~\ref{fig:correlation-heatmap}). A naive
    linear model regressing price on all 24 variables would be dominated by
    multicollinearity. A belief network instead asks which variables
    remain informative about the price once its other direct causes are
    already known, which is exactly the conditional-independence structure a
    DAG is built to represent.
  \item \textbf{Structure learning requires little prior expert knowledge.}
    Because the graph itself (not just its parameters) is learned from data,
    the model does not require the analyst to pre-specify which of the 24
    series matter, or how; it only requires selecting a reasonable candidate
    set and a learning algorithm (Section~\ref{sec:structure-learning}). This
    is valuable in a market where the dominant drivers of price are known to
    shift across regimes (demand shocks vs.\ supply shocks vs.\ pure
    financial sentiment), since a hand-specified structure risks encoding
    yesterday's regime.
  \item \textbf{Inference and interpretation are transparent.} Once learned,
    a belief network supports exact probabilistic queries (e.g.\ "what is the
    probability that WTI falls given that OPEC spare capacity has risen and
    the US dollar has strengthened?") via standard inference algorithms, and
    every edge in the graph is independently inspectable and explainable to a
    risk committee — unlike the weights of an opaque machine-learning model.
\end{enumerate}

\subsection{Advantages Over Alternative Approaches}

Relative to the two most natural alternatives for this problem, probabilistic
graphical models offer specific, concrete advantages for our use case:

\begin{itemize}
  \item \textbf{Compared to a single multivariate linear regression / VAR.} A linear
    model assumes one fixed functional form and one fixed set of
    "significant" regressors for the entire sample. Our stylized-facts
    analysis (Section~\ref{sec:stylized-facts}) shows that oil returns are
    strongly non-Normal (excess kurtosis of 18.7 for WTI) and that
    cross-variable relationships are themselves unstable over time (the
    63-day rolling correlation between WTI and the US dollar index swings
    repeatedly between $-1$ and $+1$ across the sample, Figure
    \ref{fig:rolling-corr}) — a single fixed linear structure cannot represent a relationship whose sign itself is time-varying, whereas a belief network's structure can be re-learned regime by regime.
  \item \textbf{Compared to an unconstrained machine-learning model (e.g.\ gradient boosting, a neural network).} Such models can achieve strong in-sample fit but provide little insight into why a forecast changed, which is unacceptable for risk-management sign-off. A belief network's DAG is a directly interpretable artifact: every arc corresponds to a specific, testable conditional dependence, and the Markov blanket of the price node (Section~\ref{sec:markov-chain-blanket}) tells the risk manager, in one step, the minimal set of variables that need to be monitored to screen the price node off from everything else in the system.
  \item \textbf{Combinability with regime-switching dynamics.} Because a belief network only encodes a static joint distribution, it is naturally paired with a Hidden Markov Model that first discretises each time series into a small number of latent regimes (bull/bear/stagnant), and then learns one belief network over the discretised states. 
\end{itemize}

In summary, the core modelling choice of this project — a belief network over a mixed financial/macroeconomic/geopolitical panel, informed by a preceding regime-detection step — is chosen specifically because it matches the structural features we observe in the data (Section~\ref{sec:stylized-facts}):
non-Normal, regime-dependent returns; an unstable and large set of candidate
drivers; and a need for an auditable, rather than opaque, decision process.
